\documentclass[11pt]{article}

\usepackage[margin=25mm]{geometry}

\usepackage{amsmath,amssymb,fontspec,unicode-math}

\IfFontExistsTF{DejaVu Serif}{\setmainfont{DejaVu Serif}}{\setmainfont{Latin Modern Roman}}

\IfFontExistsTF{Latin Modern Math}{\setmathfont{Latin Modern Math}}{\setmathfont{DejaVu Math TeX Gyre}}

\usepackage{xurl,hyperref}

\setlength{\parindent}{0pt}\setlength{\parskip}{6pt}

\begin{document}

\title{Discovering Opinion Intervals from Conflicts in Signed Graphs}

\author{ELVIA preliminary AI analysis}\date{21 October 2026}\maketitle

Preliminary assessment: Minor.\par

Independent scientific review remains pending. Author adoption has not been requested.

The public affinity cache can return a nonmaximal label, and the linear-time pseudocode omits exact whole-graph candidate scoring. The original conjecture was explicitly conjectural and has known author follow-up refutation.

\section*{Review scope}

Latest-version, claim-focused review of the cited author-paper passages and applicable current author-release paths. The stated counterexamples and scope controls were checked analytically; no author program, old verifier, training run or benchmark was executed.

\section*{Problem 1: A stale runner-up can produce a nonmaximal assignment}

Source: Algorithm 1, PDF p.7, line 7; pinned heuristics/src/gaic.rs, affinity/favorite updates

Assessment: supported; impact: minor

\subsection*{Source claim (paraphrase)}

Algorithm 1 requires a maximum-agreement label. In `heuristics/src/gaic.rs`, the runner-up state is updated only when the favorite changes.

\subsection*{Analysis}

After affinities become \((1,5,4)\), the cached runner-up can remain label zero. Removing two neighbors of the favorite gives \((1,3,4)\), without triggering the required rescan. Zero-temperature selection can therefore return agreement \(3\) despite an available value \(4\). This is a state-invariant failure, not a statement about every clustering run.

\subsection*{Finding and scope}

The released update can violate the maximizing assignment rule. Its frequency and benchmark impact are unknown.

\subsection*{Proposed correction}

Rescan favorites after decreases, using the existing complete affinity scan and its declared tie convention. A more selective cache update needs an equivalent maintained invariant.

An explicit local correction is described; author adoption and a corresponding published erratum are not established.

\section*{Problem 2: Exact candidate scoring has an omitted quadratic cost}

Source: Theorem 2.3, PDF p.6; Proposition B.1, PDF p.21; Algorithm 2, PDF p.22, line 11

Assessment: supported; impact: minor

\subsection*{Source claim (paraphrase)}

Theorem 2.3 states linear dependence on graph order for fixed approximation parameters; Algorithm 2 compares exact full-graph candidate agreement scores.

\subsection*{Analysis}

Two candidates such as balanced two-part and all-one assignments can differ on \(\Theta(n^2)\) edge signs. Exact comparison over arbitrary signed input requires accounting for those cross edges. Counting only sampled assignment work omits the pseudocode’s full scoring step. Approximate sampled ranking can support a separately specified randomized procedure.

\subsection*{Finding and scope}

The stated exact-scoring pseudocode is not justified by the linear runtime accounting. This does not prove that no linear-time PTAS exists.

\subsection*{Proposed correction}

Include the exact scoring cost, or explicitly sample candidate scores with an error budget compatible with the approximation proof and declared random-access input model.

An explicit local correction is described; author adoption and a corresponding published erratum are not established.

\section*{Problem 3: The conjecture has known subsequent author refutation}

Source: The original three-quarter conjecture; authors’ MLG 2026 follow-up, Theorem 3

Assessment: not\_issue; impact: no\_issue\_found

\subsection*{Source claim (paraphrase)}

The original \(3/4\) guarantee is labeled a conjecture; the authors’ MLG 2026 follow-up Theorem 3 refutes it.

\subsection*{Analysis}

For \(n=256\), random independent edge signs give a finite existence argument: there are at most \((2n)^{2n}\) interval-overlap patterns, while each fixed pattern’s chance of agreement above the target is bounded by \(\exp(-\binom n2/8)\). Their product is below one. This is a probabilistic existence argument, not an enumerated optimum or new discovery.

The elementary universal guarantee is \(|E|/2\): the better of all-overlapping and pairwise-disjoint interval assignments satisfies at least half the signed edges. The author follow-up limits a larger uniform constant in its stated asymptotic setting. This contextualizes known subsequent progress; it is not a repair of a falsely proved theorem.

\subsection*{Finding and scope}

A conjecture can be disproved by later work without constituting a false proof in the original paper. Historical attribution remains with the authors.

\section*{Limits}

Preliminary AI assessment; independent human scientific adjudication remains pending.\par

An analytic counterexample or a release discrepancy does not by itself remeasure the paper’s reported experiments.\par

\section*{Source identity}

NeurIPS 2025 official proceedings PDF (current publisher bytes, checked 2026-10-21)

\url{https://proceedings.nips.cc/paper_files/paper/2025/file/fe4326696299485549b1896c4a1afe32-Paper-Conference.pdf}

Reviewed PDF SHA-256: \nolinkurl{51924bb494762b379809806d0d6d7a7a396dad3b3f2160b3b2c930b01e05817c}

The current official publisher PDF is the same work as the original case: title, author identity and publisher hash URL agree, and the freshly downloaded scientific PDF is byte-identical to the prior source. The actual download timestamp and SHA-256 are recorded for Issue \#27. Exact-title arXiv search did not return a matching later version; background references were excluded.

\end{document}
